% 1-1-motivation.tex
%  Budget: 2pp.
%  Rules: every numeral in \lr{}; cite only from references.bib;
%  one unit at a time -- see thesis/WRITING_HANDOFF.md and APPROVALS.md.

\section{انگیزه و اهمیت موضوع}

با آزادسازی و مقررات‌زدایی از بازارهای برق در دهه‌های اخیر، قیمت برق از متغیری
تحت کنترل نهادهای متمرکز به یک سیگنال اقتصادی رقابتی بدل شده است که در بورس‌های
برق و عمدتاً از طریق حراج‌های روز-پیش تعیین می‌شود. در این ساختار، تولیدکنندگان،
خرده‌فروشان و معامله‌گران به‌طور مستقیم در معرض ریسک و نوسان قیمت قرار می‌گیرند و
توان رقابت و بقای اقتصادی آنان به دقت پیش‌بینی قیمت گره خورده است
\cite{cerjan_literature_2013, jiang_review_2018}. به همین سبب، پیش‌بینی قیمت برق
به یکی از ارکان بنیادی تصمیم‌گیری در بازارهای برق تبدیل شده است.

آنچه این مسئله را از دیگر سری‌های زمانی متمایز و دشوار می‌سازد، مجموعه‌ای از
ویژگی‌های خاص برق به‌عنوان یک کالاست: فصلی‌بودن در مقیاس‌های روزانه، هفتگی و
سالانه، لزوم تعادل پیوسته میان تولید و مصرف، و وابستگی شدید به شرایط بیرونی و
محیطی. این ویژگی‌ها موجب می‌شوند سری قیمت، نوسان‌های شدید، جهش‌های ناگهانی و
ناهمگنی واریانس (\lr{heteroscedasticity}) از خود نشان دهد
\cite{jiang_review_2018}. افزون بر این، گسترش روزافزون منابع تجدیدپذیر و وابستگی
تولید آن‌ها به شرایط جوی، بر پیچیدگی و بی‌ثباتی قیمت افزوده و پیش‌بینی آن را
دشوارتر ساخته است \cite{oconnor_review_2025, cerjan_literature_2013}. در
بازارهایی با نفوذ بالای تجدیدپذیرها، حتی پدیده‌هایی مانند قیمت‌های منفی نیز به
واقعیتی متداول بدل شده‌اند که مدل‌سازی آن‌ها را چالش‌برانگیزتر می‌کند.

شکل‌گیری قیمت برق برآیند تعامل مجموعه‌ای از عوامل بنیادی است؛ از میزان تقاضا و
بار مصرفی گرفته تا ترکیب منابع تولید \lr{—} به‌ویژه سهم منابع تجدیدپذیر \lr{—} و
شرایط جوی. پژوهش‌ها نشان داده‌اند که لحاظ‌کردن این عوامل تعیین‌کننده می‌تواند دقت
مدل‌های پیش‌بینی را به‌طور چشمگیری بهبود بخشد، هرچند درباره‌ی این‌که کدام عوامل و
به چه شماری باید در نظر گرفته شوند هنوز اجماعی وجود ندارد
\cite{chai_forecasting_2024}. همین وابستگیِ چندعاملی، در کنار ماهیت پرنوسانِ
قیمت، پیش‌بینیِ دقیق را به مسئله‌ای دشوار و چندوجهی بدل می‌کند.

اهمیت اقتصادی این مسئله را نمی‌توان نادیده گرفت. پیش‌بینی دقیق قیمت روز-پیش،
مبنای مستقیم تصمیم‌گیری در مناقصه‌گذاری (\lr{bidding})، زمان‌بندی بهره‌برداری از
واحدها (\lr{scheduling}) و مدیریت ریسک است \cite{oconnor_review_2025}، و خطا در
آن مستقیماً به زیان مالی برای بازیگران بازار ترجمه می‌شود
\cite{chai_forecasting_2024}. با افزایش سهم ذخیره‌سازها و منابع انعطاف‌پذیر در
شبکه، ارزش هر واحد بهبود در دقت پیش‌بینی نیز پیوسته فزونی می‌یابد.

در پاسخ به این نیاز، در حدود ۲۵ سال گذشته حجم گسترده‌ای از پژوهش‌ها به توسعه‌ی
روش‌های پیش‌بینی قیمت اختصاص یافته است \cite{jedrzejewski_electricity_2022}؛ از مدل‌های آماری کلاسیک مبتنی بر سری
زمانی و مدل‌های علّی گرفته تا رویکردهای مبتنی بر هوش مصنوعی
\cite{aggarwal_electricity_2009}. مرور پژوهش‌های اخیر نشان می‌دهد که این حوزه
به‌سرعت به سمت روش‌های یادگیری ماشین و یادگیری عمیق حرکت کرده و به‌ویژه مدل‌های
ترکیبی و گروهی (\lr{ensemble}) در بازار روز-پیش عملکردی قدرتمند از خود نشان
داده‌اند \cite{oconnor_review_2025, chai_forecasting_2024}.

با این‌همه، رشد سریع این حوزه با یک کاستی روش‌شناختی جدی همراه بوده است: بسیاری
از الگوریتم‌های تازه بر روی مجموعه‌داده‌های خصوصی و غیرقابل‌دسترس، در بازه‌های
آزمون کوتاه و محدود به یک بازار، و بدون مقایسه با مدل‌های ساده‌ی مرجع یا آزمون
معناداری آماری تفاوت‌ها ارزیابی شده‌اند؛ چنین رویه‌ای تشخیص این‌که کدام روش
واقعاً برتر است را ناممکن می‌کند \cite{lago_forecasting_2021}. از این‌رو، گذار
به سوی پروتکل‌های ارزیابی شفاف، تکرارپذیر و متکی بر بنچمارک‌های عمومی، خود به
ضرورتی در این حوزه بدل شده است.

پژوهش حاضر با تمرکز بر بازار برق آلمان \lr{—} از بزرگ‌ترین و پرنفوذترین بازارهای برق
اروپا با سهمی چشمگیر از منابع تجدیدپذیر \lr{—} به مسئله‌ی پیش‌بینی قیمت روز-پیش
می‌پردازد؛ بازاری که پیش‌تر با رویکردهای آماری و یادگیری عمیق مورد مطالعه قرار
گرفته است \cite{poggi_electricity_2023}. نوسان بالا و رفتار پیچیده‌ی این بازار،
آن را به بستری واقع‌گرایانه و در عین حال چالش‌برانگیز برای سنجش و مقایسه‌ی
روش‌های پیش‌بینی بدل می‌کند. اتکا به این بازار در کنار به‌کارگیری یک پروتکل
ارزیابی استاندارد و عمومی، امکان مقایسه‌ی منصفانه و تکرارپذیر روش‌های گوناگون را
فراهم می‌آورد؛ موضوعی که در فصل‌های بعد به‌تفصیل بدان پرداخته خواهد شد.
